% ══════════════════════════════════════════════════════════════════════════
\chapter{Experimento E: orientación y robustez}
\label{cap:resE}

\paragraph{Hipótesis.} Las mejoras de la fusión observadas con el punto y con los
objetos en orientación nominal se mantienen cuando la orientación del objeto es
desconocida.

\paragraph{Método (§5.6).} Monte Carlo \emph{por adquisición}: en cada repetición se
sortea una rotación uniforme en SO(3) (cuaternión gaussiano normalizado), se aplica
al objeto durante toda la apertura y es la misma en los canales mono y bi (diseño
pareado). Cubo (1000 orientaciones), caja, cilindro y tubo (200 cada uno),
$\Delta\phi=45^{\circ},90^{\circ},135^{\circ},180^{\circ}$, 800 pulsos, rejilla de $2.4\times2.4\times3$\,m. Se
corrigió así la práctica del 28/09 de sortear una orientación por posición. Se
corrieron dos receptores: TM co-polar (Exp.~E) y alineado de doble polarización
(Exp.~E2, $\Delta\phi=45^{\circ},90^{\circ},135^{\circ}$, 200 orientaciones por objeto). CSV:
\texttt{expE/expE\_mc\_<objeto>.csv} (cada repetición, con su matriz de rotación),
\texttt{expE/expE\_resumen.csv} y \texttt{expE2/expE\_resumen.csv} (medianas, IC\,95\,\%
bootstrap y probabilidad de mejora).

% generado por simulations/analise_1610/tabla_E_latex.py
\begin{table}[H]
\centering\small
\caption{Monte Carlo SO(3) por adquisici\'on: mediana de (fusi\'on N $-$ mono) y probabilidad de mejora P$_{mej}$ (Rx TM co-polar). $\Delta V$ en dB; loc.\ en m. N = 1000 (cubo), 200.}
\label{tab:E_TM}
\resizebox{\textwidth}{!}{\begin{tabular}{@{}lrr|rrrrrrrrrr@{}}
\toprule
Objeto & $\Delta\phi$ & $\bar P_{bi}/\bar P_{mono}$ & \multicolumn{2}{c}{$\Delta$SNR} & \multicolumn{2}{c}{$\Delta$Contr.} & \multicolumn{2}{c}{$\Delta$PSLR} & \multicolumn{2}{c}{$\Delta V$} & \multicolumn{2}{c}{$\Delta$loc} \\
 & & [dB] & med. & P$_{mej}$ & med. & P$_{mej}$ & med. & P$_{mej}$ & med. & P$_{mej}$ & med. & P$_{mej}$ \\
\midrule
cubo & 45° & -2.6 & -1.4 & 7\,\% & +0.3 & 81\,\% & +0.00 & 49\,\% & +0.1 & 38\,\% & +0.004 & 36\,\% \\
cubo & 90° & -20.2 & -5.7 & 0\,\% & +0.3 & 78\,\% & +0.02 & 47\,\% & +0.1 & 39\,\% & -0.001 & 59\,\% \\
cubo & 135° & -3.8 & -4.4 & 13\,\% & -0.0 & 48\,\% & +0.07 & 42\,\% & +0.9 & 43\,\% & +0.049 & 29\,\% \\
cubo & 180° & -0.0 & -2.9 & 31\,\% & -2.8 & 10\,\% & +0.11 & 34\,\% & -0.4 & 53\,\% & +0.014 & 42\,\% \\
\midrule
caja & 45° & -2.4 & -1.2 & 2\,\% & +0.4 & 98\,\% & -0.06 & 71\,\% & +0.1 & 33\,\% & +0.001 & 42\,\% \\
caja & 90° & -19.8 & -5.4 & 0\,\% & +0.6 & 91\,\% & -0.00 & 52\,\% & +0.1 & 40\,\% & -0.000 & 51\,\% \\
caja & 135° & -5.2 & -4.8 & 10\,\% & -0.5 & 37\,\% & +0.06 & 46\,\% & +0.6 & 44\,\% & +0.049 & 34\,\% \\
caja & 180° & -2.9 & -2.4 & 28\,\% & -5.2 & 3\,\% & +0.19 & 34\,\% & +0.7 & 43\,\% & +0.045 & 36\,\% \\
\midrule
cilindro & 45° & -2.6 & -1.1 & 0\,\% & +0.5 & 100\,\% & -0.00 & 50\,\% & +0.2 & 25\,\% & +0.000 & 46\,\% \\
cilindro & 90° & -20.3 & -5.5 & 0\,\% & +0.6 & 99\,\% & -0.06 & 66\,\% & +0.1 & 30\,\% & -0.001 & 60\,\% \\
cilindro & 135° & -3.1 & -5.1 & 7\,\% & -1.3 & 21\,\% & +0.48 & 26\,\% & +3.5 & 22\,\% & +0.002 & 49\,\% \\
cilindro & 180° & +0.9 & -1.6 & 28\,\% & -7.2 & 0\,\% & +0.74 & 19\,\% & +4.4 & 26\,\% & +0.061 & 14\,\% \\
\midrule
tubo & 45° & -2.4 & -1.1 & 3\,\% & +0.4 & 96\,\% & -0.00 & 50\,\% & +0.2 & 11\,\% & -0.002 & 74\,\% \\
tubo & 90° & -19.5 & -5.4 & 0\,\% & +0.6 & 94\,\% & +0.02 & 46\,\% & +0.1 & 23\,\% & -0.000 & 57\,\% \\
tubo & 135° & -3.7 & -5.6 & 16\,\% & -2.0 & 28\,\% & +0.42 & 27\,\% & +2.3 & 24\,\% & +0.195 & 40\,\% \\
tubo & 180° & -0.1 & -2.6 & 26\,\% & -7.9 & 0\,\% & +0.51 & 22\,\% & +4.5 & 31\,\% & +0.020 & 42\,\% \\
\bottomrule
\end{tabular}}
\end{table}
\begin{table}[H]
\centering\small
\caption{Monte Carlo SO(3) por adquisici\'on: mediana de (fusi\'on N $-$ mono) y probabilidad de mejora P$_{mej}$ (Rx alineado, doble polarizaci\'on). $\Delta V$ en dB; loc.\ en m. N = 200.}
\label{tab:E_AL}
\resizebox{\textwidth}{!}{\begin{tabular}{@{}lrr|rrrrrrrrrr@{}}
\toprule
Objeto & $\Delta\phi$ & $\bar P_{bi}/\bar P_{mono}$ & \multicolumn{2}{c}{$\Delta$SNR} & \multicolumn{2}{c}{$\Delta$Contr.} & \multicolumn{2}{c}{$\Delta$PSLR} & \multicolumn{2}{c}{$\Delta V$} & \multicolumn{2}{c}{$\Delta$loc} \\
 & & [dB] & med. & P$_{mej}$ & med. & P$_{mej}$ & med. & P$_{mej}$ & med. & P$_{mej}$ & med. & P$_{mej}$ \\
\midrule
cubo & 45° & -0.5 & -0.5 & 25\,\% & +0.3 & 80\,\% & +0.02 & 46\,\% & +0.2 & 35\,\% & +0.003 & 41\,\% \\
cubo & 90° & -1.4 & -2.1 & 24\,\% & +1.8 & 90\,\% & -0.05 & 53\,\% & +0.8 & 33\,\% & -0.009 & 61\,\% \\
cubo & 135° & -0.5 & -2.9 & 30\,\% & +1.2 & 78\,\% & -0.06 & 59\,\% & +0.8 & 40\,\% & -0.001 & 52\,\% \\
\midrule
caja & 45° & -0.6 & -0.4 & 22\,\% & +0.5 & 94\,\% & -0.04 & 62\,\% & +0.1 & 28\,\% & +0.001 & 40\,\% \\
caja & 90° & -1.6 & -1.5 & 20\,\% & +2.6 & 98\,\% & -0.28 & 70\,\% & +0.6 & 30\,\% & -0.001 & 54\,\% \\
caja & 135° & -1.0 & -2.7 & 21\,\% & +1.0 & 84\,\% & -0.36 & 70\,\% & +0.7 & 33\,\% & +0.002 & 49\,\% \\
\midrule
cilindro & 45° & -0.5 & -0.2 & 29\,\% & +0.5 & 100\,\% & -0.02 & 57\,\% & +0.2 & 24\,\% & -0.000 & 50\,\% \\
cilindro & 90° & -1.6 & -0.9 & 30\,\% & +2.7 & 96\,\% & -0.39 & 66\,\% & +0.8 & 23\,\% & -0.005 & 66\,\% \\
cilindro & 135° & -1.6 & -1.8 & 26\,\% & +0.9 & 80\,\% & -0.59 & 76\,\% & +1.4 & 24\,\% & -0.000 & 50\,\% \\
\midrule
tubo & 45° & -0.5 & -0.3 & 36\,\% & +0.5 & 98\,\% & -0.04 & 64\,\% & +0.3 & 15\,\% & -0.001 & 59\,\% \\
tubo & 90° & -1.0 & -1.1 & 39\,\% & +2.6 & 98\,\% & -0.08 & 65\,\% & +1.0 & 10\,\% & -0.001 & 56\,\% \\
tubo & 135° & -0.0 & -1.6 & 35\,\% & +1.0 & 92\,\% & -0.07 & 58\,\% & +1.5 & 18\,\% & +0.001 & 42\,\% \\
\bottomrule
\end{tabular}}
\end{table}


\begin{figure}[H]
\centering
\includegraphics[width=\textwidth]{E_montecarlo.pdf}
\caption{Monte Carlo SO(3) con Rx TM co-polar: distribución de (fusión N $-$ mono) por
objeto y separación; porcentajes = probabilidad de mejora.}
\label{fig:E_mc}
\end{figure}

\begin{figure}[H]
\centering
\includegraphics[width=\textwidth]{E2_montecarlo_rx_alineado.pdf}
\caption{Ídem con Rx de polarización alineada (doble polarización).}
\label{fig:E2_mc}
\end{figure}

\paragraph{Resultados e interpretación.}
\begin{enumerate}[itemsep=3pt]
  \item \textbf{Potencia (control).} Con orientación aleatoria, la potencia media del canal
  bi a 180° es la misma que la mono ($-2.9$ a $+0.9$\,dB según el objeto), como concluyó el
  28/09. Con Rx TM co-polar cae $\approx20$\,dB a 90° y 2.4--5.2\,dB a 45° y 135°
  ($\approx\cos^2\Delta\phi$): es el desajuste geométrico de polarización. Con el Rx alineado
  queda entre 0 y $-1.6$\,dB de la mono en todas las separaciones.
  \item \textbf{Rx TM co-polar:} a 45°--90° la fusión gana un contraste pequeño pero
  sistemático ($+0.3$ a $+0.6$\,dB, mejora en el 78--100\,\% de las adquisiciones) y pierde
  1.1--5.7\,dB de SNR (con el mismo presupuesto N; la fusión 2N recupera 3\,dB). PSLR y
  volumen no cambian (medianas $\approx0$). A 135°--180° la fusión empeora el contraste
  (hasta $-7.9$\,dB) y el volumen.
  \item \textbf{Rx alineado:} a 90° el contraste mejora $+1.8$ a $+2.7$\,dB con
  probabilidad del 90--98\,\% en los cuatro objetos; el PSLR mejora ligeramente
  (mediana $-0.05$ a $-0.39$\,dB, 53--70\,\%), la pérdida de SNR baja a $-0.9$ a
  $-2.1$\,dB (N) y la energía en la región del objeto sube (64--86\,\%). El precio es un
  lóbulo $+0.6$ a $+1.0$\,dB más voluminoso. A 135° la ganancia de contraste baja a
  $+0.9$ a $+1.2$\,dB.
  \item \textbf{La mejora se mantiene con la orientación sólo si el Rx está alineado
  en polarización y a 45°--90°.} La resolución y la localización no mejoran en ningún caso
  (medianas $\le0.01$\,m de diferencia).
\end{enumerate}

\conclexp{E}{Con orientación aleatoria del objeto, la fusión mejora el contraste de forma
sistemática sólo con un Rx de polarización alineada a 90° ($+1.8$ a $+2.7$\,dB en el 90--98\,\% de
las orientaciones); con un Rx TM co-polar la ganancia es marginal ($+0.3$ a $+0.6$\,dB) y cuesta
1--6\,dB de SNR, y en ningún caso mejoran la resolución ni la localización.}
